\documentclass[aspectratio=169]{beamer}
\usetheme[progressbar=frametitle]{metropolis}

\title{Gruppo I(strice) - RealSense}
\subtitle{Ricostruzione 3D Maschera}
\author{Andrea Bellu, Calvin You, Andrii Myzyn}
\date{\today}

\begin{document}

\begin{frame}
    \titlepage
\end{frame}

\begin{frame}
    \frametitle{Obiettivo Generale}
    \begin{itemize}
        \item Obiettivo: ricostruire in 3D una maschera acquisita da due Intel RealSense D435i
        \item Problema: ogni camera osserva solo una parte dell'oggetto e in un proprio sistema di riferimento
        \item Sfida: allineare nuvole di punti rumorose, parziali e con outlier
        \item Risultato atteso: modello 3D coerente con buon dettaglio geometrico e colore
    \end{itemize}
\end{frame}

\begin{frame}
    \frametitle{Obiettivo Tecnico}
    \begin{itemize}
        \item Stimare i parametri intrinseci $K, D$ di ciascuna camera
        \item Stimare i parametri estrinseci tra le camere ($R, t$)
        \item Trasformare la sorgente nel riferimento target con:
    \end{itemize}
    \[
        \mathbf{p}_{117} = \mathbf{R}\,\mathbf{p}_{108} + \mathbf{t}
    \]
    \begin{itemize}
        \item Rifinire l'allineamento con ICP (coarse + colored fine)
    \end{itemize}
\end{frame}


\end{document}